\section{Interfaces with the initial dummy region}
\label{sec:area_law_dummy_run_interfaces}

The initial dummy region has label $k_0-1$ modulo two in
\cite[Section~11]{OpenAI2026AreaLaw}. The following results pass from
an actual contact to the elementary perimeter assignment and then to
the label of an entire belt run. Write $N_0=N_{k_0}$ for the half-open
initial neighborhood and $\ell_k$ for the fine exponent at layer $k$.

% Source: 10-geometry.tex, prop:two-families, lines 154--177 and 299--323,
% revision adc7f1241b42e322a6451854ab7e4b4c146bf78a.
% These local interfaces do not assemble the full two-family partition.

\begin{theorem}[An interior elementary contact determines the whole dummy side]
    \label{thm:area_law_elementary_open_dummy_contact}
    \lean{TNLean.PEPS.AreaLaw.Geometry.fineLayer_elementarySide_subset_dummy_of_mem_openSegment}
    \leanok
    \uses{def:area_law_dyadic_layers,def:area_law_cell_fan}
    Fix an arbitrary common origin, finite endpoint set and width $C\ge 2$.
    Let $k\ge k_0$ and $z\in F_{k,\ell_k}$ be an actual fine-cell index.
    Choose any optional midpoint subdivision of its sides, and let
    $e_i=[a_i,b_i]$ be an elementary base of the resulting fan.
    If $x\in(a_i,b_i)\cap\overline{N_0}$, then
    $e_i\subseteq\overline{N_0}$.
    No lower bound on $k_0$ is required.
\end{theorem}
\begin{proof}\leanok
    \uses{thm:area_law_dummy_corner_elementary_side,
        thm:area_law_dummy_layer_disjoint}
    Choose a coarse cell in the finite dummy union whose closure contains
    $x$. Its interior is disjoint from the interior of the reference fine
    cell. After interchanging coordinates if necessary, the elementary side
    is vertical. Since $x$ lies strictly between its endpoints, disjointness
    forces its normal coordinate to be a boundary coordinate of the coarse
    cell. If a tangential boundary coordinate lay strictly between the
    elementary endpoints, the corresponding coarse-cell corner would lie
    in the open elementary side. The dummy-corner theorem excludes this.
    Thus the coarse-cell closure contains the entire elementary segment,
    and is itself contained in $\overline{N_0}$.
\end{proof}

\begin{theorem}[A triangle meets the dummy closure through its base]
    \label{thm:area_law_fan_dummy_contact}
    \lean{TNLean.PEPS.AreaLaw.Geometry.fineLayer_cellFanPolygon_inter_dummy_eq}
    \leanok
    \uses{def:area_law_dyadic_layers,def:area_law_cell_fan}
    For arbitrary common origin, finite endpoint set and natural width $C$,
    let $k\ge k_0$ and $z\in F_{k,\ell_k}$ be an actual fine-cell index.
    Choose an arbitrary optional midpoint mask. Every resulting fan triangle
    $P_i$, with elementary base $e_i$, satisfies
    \begin{align}
        P_i\cap\overline{N_0}=e_i\cap\overline{N_0}.
    \end{align}
    Empty and single-point contacts are included.
\end{theorem}
\begin{proof}\leanok
    \uses{thm:area_law_dummy_layer_disjoint,thm:area_law_cell_fan_cover,
        thm:area_law_fan_own_side_contact,thm:area_law_unsplit_side_endpoints}
    The open interior of the reference fine cell is disjoint from the
    dummy neighborhood and therefore from its closure. A common point
    belongs to the closed reference square, so it lies on its frontier.
    The product-interval frontier formula places it on one of the four
    whole sides. The own-side contact equality then places it on $e_i$.
    This proves one inclusion; the other follows from $e_i\subseteq P_i$.
\end{proof}

\begin{theorem}[Opposite dummy parity along an actual run interface]
    \label{thm:area_law_belt_run_dummy_interface}
    \lean{TNLean.PEPS.AreaLaw.Geometry.beltCellFanRun_dummy_interface}
    \leanok
    \uses{def:area_law_belt_fan_color,def:area_law_fan_runs}
    Fix a common origin, finite endpoint set, width $C\ge 2$ and
    initial layer $k_0\ge 50{,}000{,}000$, with arbitrary residue choices
    in every layer. Let $R$ be an actual run in an actual belt cell of
    layer $k\ge k_0$. If a closed segment with distinct endpoints lies
    in both the run region and $\overline{N_0}$, every constituent triangle
    of $R$ has label $k_0\ne k_0-1$ modulo two.
\end{theorem}
\begin{proof}\leanok
    \uses{thm:area_law_fan_dummy_contact,
        thm:area_law_elementary_open_dummy_contact,
        thm:area_law_belt_fan_dummy_color,thm:area_law_fan_run_color}
    The shared segment is infinite, whereas the collection of elementary
    endpoints in the fan is finite. Choose a shared point outside that
    collection and a constituent triangle containing it. The preceding
    contact equality puts the point on the triangle's base, away from both
    endpoints, hence in the open base segment. Open-contact continuation
    puts the whole base in the dummy closure. The local dummy assignment
    gives the opposite dummy parity, and constancy on the run transfers
    this label to every constituent triangle.
\end{proof}

The whole-run conclusion retains a shared nondegenerate segment. Two isolated
common points of disconnected unions do not supply this premise. Global
identifiers, isolated stars and recursive repairs remain further obligations.
